% ============================================================
% 第1章：研究背景与本文贡献
% ============================================================

\section{研究背景与问题定义}
电子元器件选型在硬件产品研发全生命周期中具有基础性作用。以 DC-DC 电源模块为例，数十家制造商提供了数万种型号，工程师需逐项核查输入输出电压范围、最大输出电流、开关频率、温度等级、封装约束、库存水平和生命周期状态等多维参数，同时查阅数据手册验证外围元件设计公式并评估供应链稳定性\cite{pecht2008parts}。该流程耗时数小时以上，且主要存在三方面瓶颈：工程知识检索高度碎片化、供应链风险缺乏系统化量化评估手段\cite{chen2021supplyrisk}、输出文档格式不够规范，难以直接支撑工程评审或 PLM 平台集成。

针对上述问题，硬禾科技旗下的 eZ-PLM 是面向电子元器件全生命周期管理的开放平台，涵盖超 110 万条在产物料与参考设计资源，支持实时器件查询和设计案例检索。同时，大语言模型（LLM）与智能体（Agent）技术在工程决策辅助领域展现出较强潜力\cite{xi2023rise,li2023llmagent}，其推理能力可与规则引擎形成互补。将 eZ-PLM 的结构化物料数据与 LLM/Agent 相结合，有望将传统人工选型转化为智能化、可追溯的自动化决策过程。
% [插图提醒：图1-1 系统应用场景示意图——展示工程师从自然语言需求输入到获得BOM选型清单、风险评估报告、拓扑分析报告的端到端流程，左侧为用户输入框，中间为Pipeline五阶段，右侧为三类输出报告]

\section{现有方案的不足}

现有元器件选型辅助方案大致分为两类：传统 PLM/ERP 方案具备真实物料数据但智能化程度有限；通用 LLM 问答方案具备自然语言交互能力但在数据真实性和可追溯性方面存在明显不足。

\subsection{传统 PLM/ERP 方案的局限}

以 SAP Integrated Product Development、PTC Windchill 等为代表的传统 PLM/ERP 平台，能够提供基于表单的参数过滤、库存查询和 BOM 导出等基础功能，但在面向复杂工程选型时存在以下局限。

\textbf{（1）缺乏自然语言交互能力。} 工程师需将设计需求手动映射为平台筛选字段和参数范围，难以直接以自然语言表达选型意图。例如“12 V 转 5 V、3 A 车规级降压方案”需拆解为类别、拓扑、电压、电流、温度等级等多个筛选条件并逐项填入表单，过程耗时且易因人为疏漏导致约束遗漏。

\textbf{（2）缺乏应用场景适配度评估能力。} 传统平台依据参数范围返回符合条件的器件列表，但难以判断器件与具体应用场景的匹配程度。例如，两款均满足“4.5--40 V 输入、5 V/3 A”的器件可能分别面向通信设备（低噪声）和工业驱动（宽温抗干扰），平台无法对这类场景差异进行智能评估，工程师仍需依赖数据手册和个人经验完成最终判断。

\textbf{（3）供应链风险评估能力薄弱。} 多数 PLM/ERP 平台主要展示库存、生命周期和供应商信息等静态数据，缺乏对供应商集中度\cite{pecht2002electronic}、交期波动和停产风险等多维因素的统一量化评估与分级预警能力\cite{goel2021supply}，尤其对于器件停产（EOL）风险，往往难以及时给出早期预警和主动替代推荐\cite{sandborn2013forecasting}。

\subsection{通用 LLM 方案的幻觉与不可溯源性}

以 ChatGPT、DeepSeek 等为代表的通用 LLM 在工程知识表达方面具有优势，但元器件选型属于高可靠性决策场景，直接依赖通用 LLM 仍面临“幻觉”（hallucination）与不可溯源问题\cite{wang2023survey}，主要表现为以下三种形式。

\textbf{（1）型号编造。} 模型可能生成看似合理但实际不存在的器件型号并附带虚构参数（如“TPS12345-Q1，5 V/3 A，AEC-Q100”），该类结果语言形式具有较强迷惑性但缺乏真实物料数据支撑。

\textbf{（2）参数捏造或误引。} 即使给出真实型号，模型也可能错误描述关键参数，此类错误可能导致严重选型偏差。

\textbf{（3）结果不可溯源。} 通用 LLM 的回答难以明确追溯到具体数据来源，无法说明推荐结论来自哪份数据手册或哪次 API 查询结果，难以满足工程评审对设计决策可追溯性的要求。

综上，两类方案之间存在明显的能力断层：传统 PLM/ERP 具备真实数据管理能力但缺乏智能化分析，通用 LLM 具备自然语言交互能力但缺乏事实约束与可溯源性\cite{lewis2020rag,gao2023retrieval}。如何将 eZ-PLM 的真实物料数据与 LLM 推理能力相结合，并通过工程规则引擎约束校验关键事实，是提升选型自动化、可靠性和可追溯性的关键。

\section{本文贡献}

本文面向 eZ-PLM 平台，设计并实现了一套融合 LLM 与工程规则引擎协同推理的智能元器件选型与风险评估 Agent，可端到端输出 BOM 选型清单、风险评估报告和拓扑分析报告三类标准化工程文档。主要技术贡献如下。

\textbf{（一）提出五阶段 Pipeline 架构}，将“需求解析—API 检索—混合评分—证据构建—报告生成”组织为可自动执行链路。在此基础上集成 LangChain ReAct Agent\cite{yao2022react}，支持多轮对话与工具自主编排，二者共同构成批量处理与交互式选型兼容的执行机制。

\textbf{(二) 设计了四级容错需求解析链}。将 LLM 语义解析与规则引擎确定性校验结合：LLM 完成初步解析后，由规则引擎通过关键词覆盖、电压方向校验和温度正则匹配逐级修正错误，将 mA 单位量级误差、否定句式漏判、LLM 拓扑幻觉等六个常见工程解析缺陷纳入自动修复框架。

\textbf{(三) 构建了以 eZ-PLM 参考设计数据为上下文约束的 LLM 混合评分机制}。评分引擎对候选器件执行规则评分（参数匹配与供应链稳定性），当 LLM 可用且存在关联参考设计时，系统基于真实设计案例对应用场景适配度和设计成熟度进行独立评分并按混合权重计算总分；LLM 不可用或数据缺失时自动回退至纯规则评分。

\textbf{(四) 建立了规则引擎事实验证与 LLM 工程叙述相分离的双层风险评估架构}。规则层对 13 类风险场景（无候选器件、停产器件、供应商过度集中等）进行确定性检测，生成结构化 \texttt{RiskItem}；LLM 层仅对已确认风险提供工程背景说明，不独立产生新的风险结论，所有 LLM 叙述均附带数据来源声明，保障风险结论的可追溯性。

\textbf{(五) 定义了可序列化为 JSON 的标准化 IR 数据模型体系}。涵盖 \texttt{PartIR}、\texttt{ScoreBreakdown}、\texttt{TopologyIR} 和 \texttt{RiskIR}，使三类工程文档与结构化 JSON 能被同步生成，直接服务于工程评审及 eZ-PLM 数据互通。

